\section{Introduction}
\label{intro}
%Motivation:
%随着视频生成基模的崛起，audio-driven video generation的效果也获得了巨大的提升，比如数字人领域中的机遇large-scaled video foundation model训练的模型，效果上相比之前有了很大的提升[omnihuman, fantasyTalk]。在实际中音频驱动的视频生成不仅可以应用在数字人演讲唱歌这种solo的场景视频中，在影视剧的人物视频生成上也有这很大的应用前景。但是目前的方法无论是数据构建上，还是训练方式上都更多的focus在了简单的talking head的场景上，使得在复杂场景上的驱动效果还没有足够好。这篇论文我们重点讨论的是如何使得audio-2-video的模型能够应用到更复杂的场景中，并且获得自然且丰富的人物动作和表情的驱动效果。
%为了实现这个目标我们认为模型需要同时具备两种能力：Text-conditioned 全局运动控制能力和Audio-conditioned 人物表情和local运动能力。
%目前的文生视频模型为我们提供了一个很好的基础模型的能力，跟大部分的方法类似，我们也是在这样的文生视频模型基础上进行finetuning，但是不同的是我们需要够保证模型在学习音频驱动的能力的同时，尽量保持住模型本身的caption的控制能力。为此我们同时利用了影视剧的视频数据和solo的talking head的数据。同时训练过程中，不像方法[FantasyTalk]那样，我们采用了全参数的模型训练，即可能的利用模型的capacity来学习足够丰富的人物视频生成的能力。14B的模型的推理效率很低，我们采用了同样的方法和数据训练了更小的1.3B的模型，从而能够降低使用成本和提高推理的效率。最后我们将开源我们的代码和模型共社区进行后续的使用和开发。
%开头需要放一个Global运动+Local运动的可视化的图片，明确论文的主要的贡献和motivation。
% With the rise of video generation foundation models, the effectiveness of audio-driven video generation has also greatly improved. For example, in the field of digital humans, models trained with large-scale video foundation models have shown significant improvements compared to previous methods [omnihuman, fantasyTalk]. In practice, audio-driven video generation can be applied not only to solo scenarios such as digital humans speaking and singing, but also has great potential in generating character videos for film and television. However, current methods, both in terms of data construction and training methods, focus more on simple talking head scenarios, resulting in insufficient driving effects in complex scenes. This paper focuses on how to enable audio-to-video models to be applied to more complex scenarios and achieve natural and rich character movements and expressions. To achieve this goal, we believe that the model needs to have two capabilities simultaneously: Text-conditioned global motion control and audio-conditioned character expressions and local motion capabilities. Current text-to-video models provide us with a good foundation model capability. Like most methods, we also fine-tune on top of such a text-to-video model, but the difference is that we need to ensure that the model can maintain its own caption control capabilities as much as possible while learning audio-driven capabilities. To this end, we simultaneously utilize video data from film and television dramas and data from solo talking heads. At the same time, unlike the [FantasyTalk] method, we adopt full-parameter model training during the training process, making it possible to utilize the model's capacity to learn sufficiently rich character video generation capabilities. Since the inference efficiency of the 14B parameter model is very low, we used the same method and data to train a smaller 1.3B parameter model, thereby reducing usage costs and improving inference efficiency. Finally, we will open source our code and models for subsequent use and development by the community.
% Fueled by the emergence of powerful video generation foundation models, audio-driven video generation has recently witnessed a significant leap in its capabilities~\cite{wan2025,kong2024hunyuanvideo}. This is particularly evident in the realm of digital humans, where models leveraging large-scale video foundation models have demonstrated substantial improvements over their predecessors~\cite{ominihuman,tian2024emo}. The practical applications of audio-driven video generation extend beyond simple scenarios like digital humans speaking or singing; it holds immense potential for generating character videos in film and television productions. However, existing methodologies often prioritize simple talking-head scenarios in both data construction and training, leading to limitations in driving more complex and dynamic scenes. 
% Some methods often generate talking heads and co-speech gestures separately, leading to less coherent outputs\cite{emo2}. Moreover, gestures produced by these methods often appear overly smooth or subdued, and lacking in diversity.
Audio-driven human video generation has made significant progress recently, largely thanks to the development and application of Diffusion models~\cite{ddpm}. Beginning with UNet-based text-to-image models and progressing to the latest DiT-based text-to-video models~\cite{wan2025,hunyuanvideo}, the quality of visual generation has also dramatically improved. Consequently, audio-driven models leveraging these latest DiT-based video foundation models are garnering increasing research attention~\cite{ominihuman,hu2025HunyuanVideo-Avatar,Wang2025FantasyTalkingRT}. However, influenced by prior work, current research primarily confines audio-driven models to single-scene human video generation, or even solely to single-character video driving. Nevertheless, in more complex scenarios such as film and television productions or multi-person scenes, audio-driven models still face numerous challenges. For instance, key questions arise: How can audio control a character while ensuring their movements are consistent and coordinated with the overall scene? How can person interactions be managed by audio and prompt jointly?
This paper primarily focuses on audio-driven human generation in such complex scenarios as film and television, aiming to enhance the efficacy of audio-driven generation through comprehensive data acquisition, robust model training, and clever yet effective inference strategies.

%介绍文字控制的重要性以及如何保持文字的控制力和音频的控制力：数据的标注.
Achieving film-quality audio-driven video, we contend, requires simultaneously leveraging the distinct yet complementary capabilities of text and audio. From a practical user perspective, text is optimally utilized for delineating the overarching dynamics of the video, including cinematic camera movements, comprehensive character trajectories, and interactions between entities. Audio, conversely, excels at dictating minute details such as character expressions and localized actions, including precise hand gestures and head orientation.
Firstly, we construct our audio-driven model by leveraging the latest Wan text-to-video foundation model\cite{wan2025}. Our aim is to integrate audio-driven capabilities while preserving its inherent text control. Crucially, to ensure our model maintains text-control fidelity during training, we utilized Qwen-VL's~\cite{Qwen2.5-VL} video understanding capabilities for detailed textual captioning of videos, with a particular emphasis on descriptions pertinent to character motion.
To effectively support generation in complex scenarios, such as film and television productions, we curated film and television-related audio-visual data from existing open-source datasets and augmented it with our own internally collected dataset of talking and singing character videos to form our comprehensive training dataset.
While some existing methods attempt to reduce training complexity by training only partial network parameters, this often leads to conflicts between text and audio control. We hypothesize that a larger model capacity is more conducive to learning superior and harmonious text and audio control, thereby mitigating such conflicts. To facilitate large-scale, full-parameter training, and drawing inspiration from established parallel training paradigms for video foundation models, we implemented a hybrid training strategy combining FSDP~\cite{fsdp} with Context Parallel, significantly accelerating the training process.
Furthermore, to ensure enhanced stability and performance, we employed a multi-stage training regimen. This includes pre-training of the audio processing modules, followed by a comprehensive pre-training phase on the entire dataset, and subsequent fine-tuning on high-quality data. Collectively, these systematic strategies enable us to develop a robust and efficient audio-driven human video generation model.


% Building upon the advancements in text-to-video models, which offer a robust foundation for video generation, our approach involves fine-tuning such a model to incorporate audio-driven capabilities. However, unlike existing methods, we prioritize maintaining the model's inherent text control capabilities while simultaneously learning to generate audio-driven content. To accomplish this, we leverage a diverse dataset comprising video data from both film and television dramas, as well as solo talking-head footage. Furthermore, in contrast to some approaches, we employ full-parameter model training to fully utilize the model's capacity and enable it to learn richer and more nuanced character video generation. 

% Recognizing the computational demands of large models, we train a smaller 1.3B parameter model, alongside the 14B parameter model, using the same data and methodology. This allows us to reduce inference costs and improve efficiency without sacrificing the quality of the generated videos. To foster further research and development in this field, we are committed to open-sourcing our code and models for the benefit of the community.
% This paper addresses the challenge of extending audio-to-video models to more complex scenarios, enabling the creation of natural and expressive character movements and expressions. To achieve this, we argue that the model must possess two essential capabilities: first, text-conditioned global motion control for directing the global dynamics of characters and the scene, and second, audio-conditioned character expressions and local motion capabilities for fine-grained, realistic performance. Building upon the advancements in text-to-video models, which offer a robust foundation for video generation, our approach involves fine-tuning such a model to incorporate audio-driven capabilities. However, unlike existing methods, we prioritize maintaining the model's inherent text control capabilities while simultaneously learning to generate audio-driven content. To accomplish this, we leverage a diverse dataset comprising video data from both film and television dramas, as well as solo talking-head footage. Furthermore, in contrast to some approaches, we employ full-parameter model training to fully utilize the model's capacity and enable it to learn richer and more nuanced character video generation. 

% Recognizing the substantial computational demands of large models, and to significantly lower the barrier to entry for users, we have trained a more lightweight model based on the Wan-1.3B architecture. This streamlined version, specifically optimized through extensive training on solo character speech and singing videos, achieves more stable and consistent results in relatively simpler scenarios.
% To foster further research and development in this field, we are committed to open-sourcing our code and models for the benefit of the community.

%介绍影视剧场景下场视频的重要性，要保持视频中运镜和场景的一致性等。为此我们提出利用FramePack类似的方式更好的解决历史视频压缩的问题，从而能够保证长视频的场景和运镜等的一致性。
Long video generation is crucial for generating videos in film and television scenarios. However, it faces challenges in maintaining stable details and consistency in scenes and even motion. Audio-driven methods like~\cite{tian2024emo} have attempted to use Motion Frames to maintain consistency between multiple clips, but an excessive number of motion frames can drastically increase computational complexity. This leads to a relatively limited number of motion frames, making it difficult to maintain long-term video stability in film and television scenarios. To address this, we introduce a approach similar to ~\cite{zhang2025framepack}, which effectively reduces the token count of Motion Frames by employing different token compression ratios at different times. This ensures the incorporation of more Motion Frames, thereby enabling the generation of more stable long videos.

% Current video foundation models possess a large number of parameters, leading to very low inference efficiency. For example, inferring 81 frames of 1080P video using a Wan14B model requires over 30 minutes on a single A100 GPU, which significantly hinders practical application and deployment. To reduce inference costs and better enable community users to create content with our models at a lower cost, we trained not only an audio-driven character video model based on Wan14B, but also a smaller audio-driven model based on Wan1.3B. By providing a smaller, more efficient model, we aim to make high-quality audio-driven video generation more accessible to a wider range of users and applications.

To train our model, we constructed a dataset containing over clips, based on both publicly available video datasets and our own collected video data. This comprehensive dataset includes videos from solo scenarios focusing on human speech and singing, as well as complex character videos from film and television dramas. 
% Through multi-stage model training, our model balances the quality and stability of generated videos with the diversity of output content. 
% Finally, we explored several potential applications of audio-driven character video generation models, such as multi-person conversation and video lip-sync editing. By controlling the audio injection area during testing and inference, we achieved control over multi-person conversations. Additionally, by adding a certain proportion of noise to the baseplate video, we achieved a character video editing function that preserves the character's original movements while driving only facial expressions and lip movements with audio, expanding the range of creative possibilities for content creators. 

Our main contributions are as follows:
\begin{itemize}
    \item \textbf{Extending Audio-Driven Generation to Complex Scenarios:} We go beyond talking heads by enabling the creation of natural and expressive character movements in diverse and challenging scenes, incorporating both text-guided global motion control and audio-driven fine-grained local motion.
    \item \textbf{Long Video Stabilization and Efficient Model Variants:} We tackle the challenges of long video generation through optimized motion frame token reduction.
    \item \textbf{Comprehensive Training Data}We leverage a large-scale, diverse dataset to train our model and validate the effectiveness of our model through extensive experiments.
\end{itemize}







